% Part: lambda-calculus
% Chapter: lambda-definability
% Section: lists

\documentclass[../../../include/open-logic-section]{subfiles}

\begin{document}

\olfileid{lam}{ldf}{lst}
\olsection{జాబితాలు}

జాబితా $[M_1, M_2, \ldots, M_n]$ను ఇలా నిర్వచించవచ్చు:
\[
  [M_1, M_2, \ldots, M_n] \ident \lambd[sz][s M_1 (s M_2 (\ldots (s M_n z)))]
\]
అనౌపచారికంగా చెప్పాలంటే, జాబితాను సూచించే పదం ఆ జాబితా యొక్క ``సంచయక'' ప్రమేయం. అది రెండు పదాలను స్వీకరిస్తుంది: మొదటిది అప్పటివరకు సంచయించిన విలువతో పాటు కొత్త మూలకాన్ని తీసుకుని కొత్త సంచిత విలువను ఇచ్చే ప్రమేయం; రెండోది ప్రారంభ సంచిత విలువ. ఈ పద్ధతి మూలకాలను ఒక్కొక్కటిగా సంచయించి చివరికి ఫలితాన్ని ఇస్తుంది.

\begin{digress}
  ప్రమేయాత్మక ప్రోగ్రామింగ్ మీకు పరిచయమైతే, ఈ నిర్వచనం \emph{కుడివైపు మడత}ను పోలి ఉంటుందని గమనిస్తారు: జాబితాను దానిలోని మూలకాల కుడివైపు మడత ప్రమేయంగా నిర్వచిస్తాం.
\end{digress}

ఈ సంకేతీకరణ ఆధారంగా కొన్ని ఉపయోగకరమైన ప్రమేయాలను సులభంగా నిర్వచించవచ్చు:
\begin{align*}
  \fn{Sum} & \ident
  \lambd[l][l \, (\lambd[xa][\fn{Add} \, x \, a])\,  \num{0}]\\
  \fn{Len} & \ident
  \lambd[l][l \, (\lambd[xa][\fn{Add} \, a \, \num{1}]) \num{0}]
\end{align*}

$\fn{Sum}$ చర్చి సంఖ్యల జాబితా మొత్తాన్ని లెక్కిస్తుంది. జాబితాపై సంచయం నిర్వహించడం ద్వారా ఇది పనిచేస్తుంది: ప్రారంభ విలువ $\num{0}$; ప్రతి మూలకానికి కొత్త విలువగా $\fn{Add} \, x\, a$ను ఇస్తుంది (ఇక్కడ $x$ ప్రస్తుత మూలకం, $a$ అప్పటివరకు సంచిత విలువ). ఫలితం అన్ని మూలకాల మొత్తం.

\begin{prob}
  $\fn{Len}$ ఏమి చేస్తుంది? వివరించండి.
\end{prob}

\begin{prob}
  జాబితాను తిరగదోసే ప్రమేయాన్ని నిర్వచించండి.
\end{prob}
\end{document}
